% !TeX root = ../main.tex
\section{Abstract}
This work investigates how embedded computer vision models for soccer analysis can be compressed and deployed efficiently on edge hardware, with a focus on object detection and pitch keypoint localisation using YOLO11 and YOLO26 model families, and using an NVIDIA Jetson Orin Nano Dev Kit as the targetted edge device. In total, \textbf{488} model artefacts were generated, and hardware acceleration, quantisation, pruning, knowledge distillation, and combinations of these techniques were explored. Alongside this, a full-stack soccer analysis system was proposed, generating high-level statistics should as formation maps, possession estimation, and pass detection.


 Results show that hardware acceleration preserved analytical performance almost exactly while reducing latency by \textbf{31.3\%} for object detection and \textbf{29.9\%} for pose estimation, with mean mAP$_{50\text{--}95}$ changes of only \textbf{-0.3\%} and \textbf{-0.2\%} respectively. Quantisation produced the largest latency gains at \textbf{56.6\%} for object models and \textbf{59.6\%} for pose models, but with larger mean mAP$_{50\text{--}95}$ losses of \textbf{4.9\%} and \textbf{6.5\%}. Pruning delivered moderate latency gains, with a relatively small accuracy cost for object detection but a severe penalty for pose estimation, showing that the keypoint models were far less pruning-tolerant. Knowledge distillation did not provide a direct latency benefit in the checkpoint-level comparison and produced smaller accuracy shifts overall, making it more useful as a supporting technique than as a direct compression method.

From a hardware perspective, quantisation delivered the strongest system-level efficiency gains overall, reducing average power draw by \textbf{42.3\%} and \textbf{41.2\%} and average temperature by \textbf{13.1\%} and \textbf{14.0\%} for object and pose models respectively, although it also increased average CPU usage by about \textbf{13\%} across both tasks.


Under the strongest replay setting, the best full-stack configuration achieved a Pass mAP$_{1\text{--}5}$ of 0.697 and 39.61\% GPU utilisation, while reaching an effective throughput of 13.02 FPS on the target device. Overall, the thesis shows that deployment-focused compression can substantially improve the practicality of embedded soccer analytics while retaining useful analytical performance, and that a full computer vision pipeline for match statistic generation can be run effectively on constrained edge hardware.
